\title{Controlling Text-to-Image Diffusion by Orthogonal Finetuning}

\begin{abstract}
Text-to-image diffusion models have recently demonstrated remarkable abilities in synthesizing photorealistic visuals from textual input. However, steering these sophisticated models for various downstream applications presents a major research challenge. To address this, we propose Orthogonal Finetuning~(OFT), a systematic approach for tailoring text-to-image diffusion models to new tasks. OFT differentiates itself from prior approaches by provably maintaining hyperspherical energy, a property representing the inter-neuron angular relationships on the unit hypersphere. Our analysis shows this invariance is essential for retaining the semantic content generation performance of these models. For added robustness during adaptation, we develop Constrained Orthogonal Finetuning (COFT), which integrates a radius constraint alongside the hyperspherical condition. Our investigation centers on two key finetuning applications: (1) subject-driven generation, wherein the objective is to synthesize subject-specific imagery from a limited set of example images and accompanying textual guidance; and (2) controllable generation, which involves leveraging supplemental control inputs. Through empirical studies, we demonstrate that the OFT framework achieves superior image quality and faster convergence compared to state-of-the-art baselines.
\end{abstract}

\section{Introduction}

The advent of text-to-image diffusion models~\cite{saharia2022photorealistic,ramesh2022hierarchical,rombach2022high} has set a new benchmark for generating realistic images with textual instructions. Nonetheless, using only text prompts may fall short when it comes to specifying detailed or precise aspects of the output. To mitigate this limitation, we focus on two representative categories of text-to-image generation challenges in this work:

\begin{itemize}[leftmargin=*,nosep]

    \item \textbf{Subject-driven generation}~\cite{ruiz2023dreambooth}: Provided a handful of exemplar images depicting a particular subject, this task entails creating new images that depict the same subject within different scenarios, guided by textual prompts.
    \item \textbf{Controllable generation}~\cite{zhang2023adding,mou2023t2i}: Here, additional cues such as canny edge maps or segmentation labels supplement the text prompt, with the goal of generating images that adhere to both forms of guidance.
\end{itemize}

At their core, both tasks address the challenge of refining text-to-image diffusion models in a way that retains their original generative capabilities acquired during pretraining. We outline the key requirements for a robust finetuning approach as follows: (1) \emph{training efficiency}, which involves minimizing both the number of trainable parameters and the training epochs required; and (2) \emph{generalizability preservation}, or maintaining the model's capacity for producing high-fidelity, varied outputs. Common finetuning techniques include adjusting the model’s neuron weights with a minimal learning rate (\emph{e.g.}, \cite{ruiz2023dreambooth}) or integrating a lightweight, re-parameterized component to the network (\emph{e.g.}, \cite{hulora2022,zhang2023adding}). Although these approaches are straightforward, they do not provide assurances that the generative performance from pretraining is retained. Furthermore, there is a lack of systematic guidance on devising optimal finetuning strategies or selecting appropriate hyperparameters such as training duration. Central to this challenge is the absence of a metric for assessing how much of the pretrained generative capability is preserved. Present methods generally rely on the implicit belief that minimizing the Euclidean distance between finetuned and original model parameters equates to better retention of prior abilities. As a consequence, finetuning tends to use very small learning rates. However, while sometimes valid, we contend that Euclidean distance alone falls short as an indicator of semantic preservation. Thus, employing a more nuanced, structural metric to distinguish between the finetuned and pretrained models could significantly enhance both the stability of finetuning and the retention of pretraining performance.

Motivated by empirical findings that hyperspherical similarity effectively encapsulates semantic relationships~\cite{liu2017deep,liu2018decoupled,chen2020angular}, we utilize hyperspherical energy~\cite{liu2018learning} to analyze the pairwise relational patterns among neurons. Specifically, hyperspherical energy is computed as the aggregate hyperspherical similarity (for instance, cosine similarity) between every pair of neurons within a layer, thereby quantifying how uniformly neurons are distributed over the unit hypersphere~\cite{liu2021learning}. We propose that optimal finetuning should yield minimal alteration in hyperspherical energy relative to the original pretrained model. One straightforward approach is to incorporate a regularization term that maintains hyperspherical energy throughout finetuning; however, this strategy does not ensure effective minimization of the energy difference. To address this, we leverage the invariance property of hyperspherical energy: when an identical orthogonal transformation is applied to all neurons, the pairwise hyperspherical similarity remains invariant. Drawing on this insight, we introduce Orthogonal Finetuning~(OFT), a method designed to adapt large text-to-image diffusion models for downstream applications without altering their hyperspherical energy. The foundational concept of OFT involves learning an orthogonal transformation—shared across the layer—that preserves the inter-neuronal angles. This approach can also be interpreted as a redefinition of the standard coordinate framework for neurons within each layer. Additionally, recognizing that maintaining a smaller Euclidean distance between the finetuned and pretrained model is beneficial for retaining pretrained capabilities, we propose a variant called Constrained Orthogonal Finetuning~(COFT). This extension restricts the finetuned parameters to reside within a hypersphere of fixed radius centered at the pretrained neuron locations.

The rationale behind using orthogonal transformations for finetuning neurons draws in part from the analogy to the 2D Fourier transform, a process by which an image is separated into its magnitude and phase spectra. Critically, the phase spectrum encapsulates angular information between the input and its basis and is primarily responsible for conveying the image's semantic content. Notably, one can reconstruct the original image’s semantics by combining its phase spectrum with an arbitrary magnitude spectrum. This insight points to the importance of modifying neuron directions to alter semantics in generated images, a fundamental objective in both subject-driven and controlled image generation. Nonetheless, unconstrained alterations in neuron orientations can undermine the generative capability instilled during pretraining. To mitigate this risk, we advocate for maintaining the angles between every pair of neurons, drawing on the premise that these angular relationships encapsulate core network knowledge. Accordingly, we propose learning a layer-wise, shared orthogonal transformation for neurons so that the hyperspherical energy remains constant.

Our approach is further informed by the framework of orthogonal over-parameterized training~\cite{liu2021orthogonal}, wherein classification networks are trained from randomly initialized states by subjecting their neurons to orthogonal transformations. This strategy leverages the fact that a randomly initialized network is expected to exhibit low hyperspherical energy, and the central goal in~\cite{liu2021orthogonal} is to maintain this low energy throughout training (as reduced energy correlates with improved generalization in classification tasks~\cite{liu2018learning,lin2020regularizing}). The findings in~\cite{liu2021orthogonal} demonstrate that orthogonal transformations possess ample expressiveness to successfully train networks for classification. In contrast, our focus lies on finetuning text-to-image diffusion models with the objectives of enhanced controllability and superior generative capabilities on downstream tasks. We stress two distinctions between our OFT method and the approach in~\cite{liu2021orthogonal}. Firstly, while~\cite{liu2021orthogonal} pursues minimization of hyperspherical energy, OFT instead enforces its preservation relative to the pretrained model, ensuring that the underlying pretrained structure is not disrupted by subsequent finetuning. In the context of diffusion models, reducing hyperspherical energy may be detrimental to original semantic representations. Secondly, OFT explicitly minimizes deviations from the pretrained model, leading to a constrained optimization variant, whereas~\cite{liu2021orthogonal} does not impose such restrictions. Finetuning thus involves balancing adaptability with retention of learned features, a balance that our OFT framework is designed to optimize. Our contributions are summarized as follows:

\begin{itemize}[leftmargin=*,nosep]

    \item We introduce Orthogonal Finetuning (OFT), an innovative approach for finetuning text-to-image diffusion models with the aim of enhancing model controllability. To address training stability, we develop a constrained version of this method that restricts angular deviations from the original pretrained parameters.
    \item Distinct from standard finetuning strategies, OFT conducts model adaptation while mathematically guaranteeing the conservation of hyperspherical energy—an empirical indicator strongly correlated with maintaining the pretrained model’s generative semantics.
    \item We evaluate OFT on subject-driven and controllable generation tasks, carrying out an extensive set of experiments that establish superior performance over existing approaches. Notably, OFT yields gains in image generation quality, convergence rates, and stability throughout the finetuning process. In addition, OFT is sample-efficient, attaining effective convergence with markedly fewer images and training epochs.
    \item For tasks requiring controllable generation, we present a novel control consistency metric. This metric involves extracting the control signal from each generated sample and quantitatively comparing it to the original input control, offering additional evidence for OFT's robust controllability.
\end{itemize}

\section{Related Work}

\textbf{Text-to-image diffusion models}. Recent years have witnessed significant breakthroughs in text-to-image synthesis~\cite{saharia2022photorealistic,ramesh2022hierarchical,rombach2022high,gu2022vector,nichol2022glide}, largely propelled by advances in diffusion-based generation~\cite{ho2020denoising,song2021score,song2021denoising,dhariwal2021diffusion} and multimodal vision-language representation learning~\cite{su2020vlbert,lu2019vilbert,radford2021learning,wang2021simvlm,li2022blip,li2023blip,alayrac2022flamingo,schuhmann2022laion}. Approaches such as GLIDE~\cite{nichol2022glide} and Imagen~\cite{saharia2022photorealistic} implement diffusion in pixel space. While GLIDE integrates the training of both text encoder and diffusion model using paired datasets, Imagen operates with a fixed pretrained text encoder. Alternatively, Stable Diffusion~\cite{rombach2022high} and DALL-E2~\cite{ramesh2022hierarchical} perform diffusion within a latent space; the former leverages VQ-GAN~\cite{esser2021taming} for latent visual codebook construction, and the latter utilizes CLIP~\cite{radford2021learning} for multimodal joint embedding. Beyond diffusion frameworks, research has also explored generative adversarial networks~\cite{reed2016generative,zhang2017stackgan,xu2018attngan,li2019controllable} and autoregressive models~\cite{ramesh2021zero,ding2021cogview,wu2022nuwa,yuscaling} for text-conditioned image generation. OFT remains methodologically agnostic, enabling its application to any diffusion-based text-to-image model.

\textbf{Subject-driven generation}. In efforts to prevent unwanted changes to target subjects, works such as~\cite{avrahami2022blended,nichol2022glide} integrate user-supplied masks as auxiliary inputs. Inversion-based approaches~\cite{choi2021ilvr,dhariwal2021diffusion,ramesh2022hierarchical,gal2022image} allow for editing contextual attributes while keeping the subject fixed. Editing without explicit mask guidance is enabled by~\cite{hertz2022prompt}, supporting both local and global changes. These strategies, however, often fail to robustly retain subject identity. Pivotal Tuning~\cite{roich2022pivotal} adjusts the generator around an initially inverted code with explicit identity regularization, and~\cite{nitzan2022mystyle} learns individualized generative priors from curated facial datasets. \cite{casanova2021instance} achieves instance-level variation, yet at the expense of detail retention. DreamBooth~\cite{ruiz2023dreambooth} customizes models for novel subjects by jointly optimizing reconstruction and class-specific preservation losses, given a unique token and limited examples. The OFT method utilizes DreamBooth's setup, but diverges by finetuning via orthogonal transformations instead of simply applying gradient descent with a reduced learning rate.

\textbf{Controllable generation}. Image-to-image translation is essentially a type of controllable generation, and has traditionally relied on conditional generative adversarial networks~\cite{isola2017image,zhu2017toward,wang2018high,park2019semantic,choi2018stargan}. Recent advancements have seen the use of diffusion models in this domain~\cite{saharia2022palette,voynov2022sketch,wang2022pretraining}. Notably, ControlNet~\cite{zhang2023adding} introduces a strategy to steer pretrained diffusion models via finetuning, incorporating additional control cues and achieving state-of-the-art results in controllability. Similarly, T2I-Adapter~\cite{mou2023t2i}, developed concurrently, also adapts pretrained diffusion models for enhanced control over image generation. In alignment with the evaluation protocol from~\cite{zhang2023adding,mou2023t2i}, we leverage OFT for finetuning pretrained diffusion models, consistently attaining superior controllability while requiring less training data and fewer parameters to be updated. Importantly, OFT imposes no extra computational cost at inference.

\textbf{Model finetuning}. The practice of adapting large pretrained models to downstream applications through finetuning has garnered increasing attention in recent years~\cite{devlin2018bert,brown2020language,he2022masked}. Adaptation techniques, such as those discussed in \cite{hulora2022,houlsby2019parameter,pfeiffer2020adapterfusion}, are extensively explored, especially in the context of natural language processing. Of particular relevance is LoRA~\cite{hulora2022}, which posits a low-rank approach for additive weight modification during finetuning. Diverging from this method, OFT adopts a layer-wise shared orthogonal transformation, updating neural weights multiplicatively. This approach preserves the inter-neuronal angular relationships within a layer, resulting in improved stability.

\section{Orthogonal Finetuning}

\subsection{Why Does Orthogonal Transformation Make Sense?}

We begin by examining the motivations for employing orthogonal transformation in the finetuning of text-to-image diffusion models. To clarify this issue, we split it into two fundamental aspects: (1) the rationale behind finetuning neuron angles (i.e., their directional orientation), and (2) the justification for utilizing orthogonal transformations specifically to adjust these angles.

To address the first question, we build upon the empirical findings of \cite{liu2018decoupled,chen2020angular}, which indicate that differences in angular features serve as effective descriptors of the semantic gap. SphereNet~\cite{liu2017deep} demonstrated that constraining all neurons of a neural network onto a unit hypersphere not only maintains the representational capacity but also improves generalization, suggesting that the informational essence of data can be predominantly represented by neuron directions. To further elucidate the critical role of neuron orientation, we present a toy experiment in Figure~\ref{fig:toy_ae}. In this experiment, a conventional convolutional autoencoder is trained on images of flowers, using the standard inner product for feature map computation during training ($z$ is the convolution output for kernel $\bm{w}$ and input $\bm{x}$ within the receptive field). At the inference phase, we investigate three different approaches for deriving feature maps: (a) replicating the inner product calculation used during training, (b) isolating magnitude-based information, and (c) focusing on angular information. The findings illustrated in Figure~\ref{fig:toy_ae} demonstrate that the angular components alone nearly reconstruct the original images with high fidelity, whereas the magnitude information fails to provide meaningful content. It is important to note that the cosine-based activation in (c) is not involved during training—the training strictly utilizes the inner product. These observations underscore that neuron orientation predominantly encodes semantic content within the inputs, implying that adjustments to neuron directions could be more beneficial for altering semantic attributes in images.

Turning to the second question, a direct approach for adapting neuron orientations is to synchronously apply rotations or reflections to all neurons within a layer, an operation naturally associated with orthogonal transformations. While more intricate angular transformations, such as those assigning distinct rotation angles to different neurons, offer added flexibility, we observe that orthogonal transformations strike an optimal balance between adaptability and structured regularization. Furthermore, \cite{liu2021orthogonal} establishes that orthogonal transformations suffice for effective neural network training. To empirically substantiate the benefits of imposing orthogonality, we conduct a controllable generation experiment within the ControlNet~\cite{zhang2023adding} context; results are displayed in Figure~\ref{fig:ortho}. Our experiments show that the conventional OFT maintains reliable performance and allows for precise control after 20 training epochs. Conversely, removing the orthogonality constraint from OFT results in the inability to synthesize realistic images or exert meaningful control. This experiment confirms the crucial regularization effect introduced by the orthogonality requirement in OFT.

\subsection{General Framework}

Typically, conventional finetuning involves the application of gradient descent with a small step size to adjust the parameters of a neural network (either globally or in selected layers). Using a reduced learning rate implicitly restricts the model from deviating significantly from the initialization, with traditional finetuning aiming to adjust weights such that $\thickmuskip=2mu \medmuskip=2mu \| \bm{M} - \bm{M}^0\|$ remains small; here, $\bm{M}$ and $\bm{M}^0$ refer to the current and pretrained model parameters, respectively. Although this implicit proximity is sensible, it may be overly permissive, particularly for the adaptation of large models. To provide further restriction, LoRA introduces an explicit constraint on the weight change, demanding that $\thickmuskip=2mu \medmuskip=2mu \text{rank}(\bm{M} - \bm{M}^0)=r'$, where $r'$ is deliberately chosen to be low. In contrast, OFT operates by imposing a condition on the inter-neuron relationships, specifically enforcing $\thickmuskip=2mu \medmuskip=2mu \| \text{HE}(\bm{M})-\text{HE}(\bm{M}^0) \|=0$, where $\text{HE}(\cdot)$ indicates the hyperspherical energy for a set of weights. To illustrate this approach, consider the fully connected layer given by $\thickmuskip=2mu \medmuskip=2mu \bm{W}=\{\bm{w}_1,\cdots,\bm{w}_n\}\in\mathbb{R}^{d\times n}$, where the $i$-th neuron is represented as $\bm{w}_i\in\mathbb{R}^{d}$ (with $\bm{W}^0$ being the layer’s pretrained weights). The output of this layer is $\thickmuskip=2mu \medmuskip=2mu \bm{z}\in\mathbb{R}^n$ calculated by $\thickmuskip=2mu \medmuskip=2mu \bm{z}=\bm{W}^\top\bm{x}$, given an input $\thickmuskip=2mu \medmuskip=2mu \bm{x}\in\mathbb{R}^d$. In this setup, OFT aims to minimize the change in hyperspherical energy resulting from finetuning relative to the pretrained state, formally expressed as

\begin{equation}\label{eq:oft_principle}
\footnotesize
\min_{\bm{W}} \left\| \text{HE}(\bm{W})-\text{HE}(\bm{W}^0)\right\|~~\Leftrightarrow~~\min_{\bm{W}} \bigg{\|} \sum_{i\neq j} \|\hat{\bm{w}}_i-\hat{\bm{w}}_j\|^{-1}-\sum_{i\neq j} \|\hat{\bm{w}}^0_i-\hat{\bm{w}}^0_j\|^{-1}\bigg{\|}
\end{equation}

where the normalized neurons are given by $\thickmuskip=2mu \medmuskip=2mu \hat{\bm{w}}_i:=\bm{w}_i/\|\bm{w}_i\|$, and hyperspherical energy for weights $\bm{W}$ is defined by $\thickmuskip=2mu \medmuskip=2mu \text{HE}(\bm{W}):= \sum_{i\neq j} \|\hat{\bm{w}}_i-\hat{\bm{w}}_j\|^{-1}$. Notice that the minimum value attainable for Eq.~\eqref{eq:oft_principle} is zero. This minimum is realized whenever $\bm{W}$ is related to $\bm{W}^0$ by an orthogonal transformation (either rotation or reflection), i.e., $\thickmuskip=2mu \medmuskip=2mu \bm{W}=\bm{R}\bm{W}^0$, with $\thickmuskip=2mu \medmuskip=2mu \bm{R}\in\mathbb{R}^{d\times d}$ being an orthogonal matrix (where $\det(\bm{R})$ is $1$ for rotations and $-1$ for reflections). This principle underlies OFT: only orthogonal, layer-wide

\begin{equation}\label{eq:oft_general}
    \footnotesize
    \bm{z}=\bm{W}^\top\bm{x}=(\bm{R}\cdot\bm{W}^0)^\top\bm{x},~~~\text{s.t.}~\bm{R}^\top\bm{R}=\bm{R}\bm{R}^\top=\bm{I}
\end{equation}

In this formulation, $\bm{W}$ is the weight matrix fine-tuned via OFT, while $\bm{I}$ represents the identity matrix. The overall mechanism of OFT is depicted in Figure~\ref{fig:block}. Drawing an analogy to LoRA's zero initialization, it is essential for OFT to perform fine-tuning strictly starting from the original parameters $\bm{W}^0$. To enforce this behavior, the orthogonal matrix $\bm{R}$ is initially set to the identity, thereby ensuring that the adapted model begins with the pretrained weights. Maintaining the orthogonality constraint on $\bm{R}$ can be accomplished through differentiable orthogonalization methods, as described by \cite{liu2021orthogonal,lezcano2019cheap}. In the following, we elaborate on strategies for achieving orthogonality efficiently during optimization.

\subsection{Efficient Orthogonal Parameterization}\label{sect:eff_ortho}

Although differentiable, conventional orthogonalization techniques such as the Gram-Schmidt process are computationally burdensome for large-scale applications~\cite{liu2021orthogonal}. To improve computational tractability, we utilize the Cayley parameterization, which generates orthogonal matrices in a more efficient manner. Specifically, we define the orthogonal matrix as $\thickmuskip=2mu \medmuskip=2mu \bm{R}=(\bm{I}+\bm{Q})(I-\bm{Q})^{-1}$, where $\bm{Q}$ is a skew-symmetric matrix satisfying $\thickmuskip=2mu \medmuskip=2mu \bm{Q}=-\bm{Q}^\top$. One caveat is that the Cayley transform produces only orthogonal matrices with determinant equal to $1$, restricting $\bm{R}$ to the special orthogonal group. However, in our empirical studies, we observe that this restriction does not hinder model performance. Nevertheless, even with the Cayley parameterization, parameter efficiency can deteriorate for large $d$, as $\bm{R}$ may become unwieldy. To circumvent this, we partition $\bm{R}$ into a block-diagonal matrix composed of $r$ blocks, leading to the following structure:

\begin{equation}
    \footnotesize
    \bm{R}=\text{diag}(\bm{R}_1,\bm{R}_2,\cdots,\bm{R}_r)=\begin{bmatrix}
    \bm{R}_1\in \text{O}(\frac{d}{r}) &  &\\
    &\ddots & \\
    & & \bm{R}_r\in \text{O}(\frac{d}{r})
    \end{bmatrix}\in \text{O}(d)
\end{equation}

Here, $\text{O}(d)$ refers to the orthogonal group in $d$ dimensions, and both $\thickmuskip=2mu \medmuskip=2mu \bm{R}\in\mathbb{R}^{d\times d}$ and each block matrix $\thickmuskip=2mu \medmuskip=2mu \bm{R}_i\in\mathbb{R}^{d/r\times d/r}$ (for all $i$) are orthogonal. Setting $\thickmuskip=2mu \medmuskip=2mu r=1$ yields the ordinary, unconstrained orthogonal matrix. For a general $d\times d$ orthogonal matrix, the number of independent parameters is $\thickmuskip=2mu \medmuskip=2mu d(d-1)/2$, leading to computational complexity $\mathcal{O}(d^2)$. When using an $r$-block diagonal orthogonal structure, there are $\thickmuskip=2mu \medmuskip=2mu d(d/r-1)/2$ parameters, corresponding to complexity $\thickmuskip=2mu \medmuskip=2mu \mathcal{O}(d^2/r)$. By enforcing block-sharing—i.e., using identical block matrices such that $\thickmuskip=2mu \medmuskip=2mu \bm{R}_i = \bm{R}_j$ for any $i\neq j$—the parameter count can be reduced to $\mathcal{O}(d^2/r^2)$. Notably, all these parameter reduction approaches maintain the orthogonality of $\bm{R}$, thereby ensuring preservation of hyperspherical energy.

We analyze OFT's parameter efficiency in comparison to LoRA. For LoRA configured with low-rank parameter $r'$, its trainable parameter count amounts to $\thickmuskip=2mu \medmuskip=2mu r'(d+n)$. Assuming both $r$ and $r'$ scale with model width $d$, such that $\thickmuskip=2mu \medmuskip=2mu r = r' = \alpha d$ where $\thickmuskip=2mu \medmuskip=2mu 0<\alpha\leq 1$ is a fixed constant, the parameter complexity for LoRA increases to $\mathcal{O}(d^2 + dn)$, while OFT achieves a lower complexity of $\mathcal{O}(d)$. To demonstrate this gap, consider a weight matrix sized $\thickmuskip=2mu \medmuskip=2mu 128\times 128$; under this configuration, LoRA introduces $2{,}048$ trainable parameters when $\thickmuskip=2mu \medmuskip=2mu r'=8$, whereas OFT requires $960$ trainable parameters for $\thickmuskip=2mu \medmuskip=2mu r=8$ (without block sharing).

\subsection{Constrained Orthogonal Finetuning}

The expressiveness of standard OFT can be further curbed by restricting the adapted model to stay close to the initial pretrained parameters. In particular, COFT modifies the forward computation as follows:

\begin{equation}\label{eq:coft_general}
    \footnotesize
    \bm{z} = \bm{W}^\top\bm{x} = (\bm{R}\cdot\bm{W}^0)^\top\bm{x},~~~\text{s.t.}~\bm{R}^\top\bm{R} = \bm{R}\bm{R}^\top = \bm{I},~\norm{\bm{R}-\bm{I}}\leq \epsilon
\end{equation}

In this formulation, both orthogonality and an $\epsilon$-bounded deviation from the identity are imposed. The orthogonality condition can be met using the Cayley parameterization (see Section~\ref{sect:eff_ortho}), although incorporating the $\epsilon$-deviation constraint into this framework remains challenging. To build intuition for the Cayley transform, we apply the Neumann series expansion to the expression $\thickmuskip=2mu \medmuskip=2mu \bm{R}=(\bm{I}+\bm{Q})(\bm{I}-\bm{Q})^{-1}$, yielding the approximation $\thickmuskip=2mu \medmuskip=2mu \bm{R}\approx\bm{I}+2\bm{Q

The convergence of the series in the operator norm permits us to transfer the constraint $\thickmuskip=2mu \medmuskip=2mu\|\bm{R}-\bm{I}\|\leq\epsilon$ directly into the Cayley transform. This yields an analogous constraint $\thickmuskip=2mu \medmuskip=2mu \|\bm{Q}-\bm{0}\|\leq\epsilon'$, where $\bm{0}$ designates a zero matrix and $\epsilon'$ is a separate error tolerance distinct from $\epsilon$. The imposition of this new condition on the matrix $\bm{Q}$ is straightforward with projected gradient descent. To initialize the orthogonal matrix $\bm{R}$ as the identity, we simply set $\bm{Q}$ to be a zero matrix. COFT can thus be understood as enforcing two explicit criteria: limited deviation in hyperspherical energy and a tightly regulated difference from the initial model parameters. Whereas most finetuning strategies typically impose the latter restriction implicitly, COFT makes it an explicit optimization target. While OFT already achieves robust results, our findings indicate that COFT improves stability in finetuning further by explicitly constraining model deviation. The influence of $\epsilon$ on COFT’s effectiveness is illustrated in Figure~\ref{fig:eps_coft}. The figure reveals that $\epsilon$ modulates the degree of adaptation possible: increasing $\epsilon$ renders COFT closer in behavior to OFT, while decreasing $\epsilon$ results in outputs that more closely resemble the original pretrained diffusion model.

\subsection{Re-scaled Orthogonal Finetuning}

We introduce a straightforward extension to standard OFT by also learning a separate scaling factor per neuron. This addition is driven by the insight that scaling neurons has no effect on the hyperspherical energy, as normalization nullifies the scaling during computation. The revised forward operation takes the form $\thickmuskip=2mu \medmuskip=2mu\bm{z}=(\bm{R}\bm{W}^0\bm{D})^\top\bm{x}$,\footnote{\textbf{Errata}: In the NeurIPS camera ready version, there was a typographical error in the re-scaled OFT forward pass, given incorrectly as $\thickmuskip=2mu \medmuskip=2mu\bm{z}=(\bm{D}\bm{R}\bm{W}^0)^\top\bm{x}$. The original code uses the correct formulation; thus, the findings in Appendix~\ref{app:rescaled} remain accurate.} where $\thickmuskip=2mu \medmuskip=2mu \bm{D}=\text{diag}(s_1,\cdots,s_n)\in\mathbb{R}^{n\times n}$ is a learnable diagonal matrix with all entries $s_1,\cdots,s_n$ strictly positive. Whereas classical OFT (as defined in Eq.~\eqref{eq:oft_general}) only treated $\bm{R}$ as a trainable parameter, the re-scaled variant introduces both $\bm{D}$ and $\bm{R}$ as learnable matrices. This generalization endows OFT with greater adaptability at the cost of a minimal parameter increase. For most of our empirical evaluations, we adhere to basic OFT to isolate the impact of orthogonal transformations; nevertheless, our observations confirm that re-scaled OFT tends to yield superior results in general (refer to Appendix~\ref{app:rescaled}).

\section{Intriguing Insights and Discussions}

\textbf{OFT is adaptable across model types}. The OFT methodology is, in principle, compatible with all neural network architectures. In the Transformer setting, techniques such as LoRA are commonly used for modifying attention weights~\cite{hulora2022}. To ensure a meaningful comparison with LoRA, our application of OFT in experiments is also restricted to the attention parameters. Furthermore, beyond dense layers, OFT naturally extends to convolutional layers; the block-diagonal pattern in $\bm{R}$ presents interesting implications for convolutions—a notable distinction from LoRA. For example, setting the block count equal to the number of input channels allows each block to exclusively transform a particular channel, analogous to independently learned depth-wise convolution filters~\cite{chollet2017xception}. Conversely, by sharing all $\bm{R}$ blocks, OFT enacts a uniform orthogonal transformation across all neuron channels. We offer initial results on applying OFT to convolutional layers in Appendix~\ref{app:convolution}.

\textbf{Connection to LoRA}. LoRA mitigates abrupt alterations to the pretrained weights by introducing a low-rank matrix, thereby preserving the original information. In contrast, OFT employs an orthogonal (and thus full-rank) transformation on the pretrained weight matrix, constraining this modification to be information-preserving. The forward propagation in OFT can be reformulated as $\thickmuskip=2mu \medmuskip=2mu \bm{z}=(\bm{R}\bm{W}^0)^\top\bm{x}=(\bm{W}^0+(\bm{R}-\bm{I})\bm{W}^0)^\top\bm{x}$, where $\thickmuskip=2mu \medmuskip=2mu (\bm{R}-\bm{I})\bm{W}^0$ is comparable to the low-rank update mechanism in LoRA. As $\bm{W}^0$ is generally full-rank, OFT can also implement a low-rank weight update whenever $\thickmuskip=2mu \medmuskip=2mu \bm{R}-\bm{I}$ is low-rank. Analogous to the rank parameter $r'$ in LoRA, OFT introduces a diagonal block size $r$ to restrict the number of parameters being optimized. Importantly, LoRA and OFT each embody unique paradigms for parameter efficiency: LoRA harnesses low-rank factorization, whereas OFT leverages sparsity, specifically block-diagonal orthogonality, to attain efficiency.

\textbf{Why OFT converges faster}. As illustrated in Figure~\ref{fig:toy_ae}, changing neuron directions is the most direct means of semantic adaptation—precisely the modification that OFT enables. Additionally, OFT can be interpreted as adjusting neurons on a smooth hyperspherical manifold, resulting in a more favorable optimization landscape, an observation that is also supported empirically by \cite{liu2021orthogonal}.

\textbf{Why not minimize hyperspherical energy}. One notable distinction from \cite{liu2021orthogonal} is that our method does not target the minimization of hyperspherical energy. In classification contexts, it is preferable for neurons to exhibit non-redundancy. Minimizing hyperspherical energy enforces uniform neuron distribution over the hypersphere, which is not aligned with finetuning goals as it could compromise valuable pretrained knowledge.

\textbf{Balancing flexibility and regularity in finetuning}. Our investigation reveals a fundamental tension between flexibility and regularity in finetuning approaches. While standard finetuning offers maximum adaptability, it often suffers from instability, making the model susceptible to collapse. Interestingly, OFT—despite its simplicity—achieves a favorable compromise between these competing factors by maintaining the angular relationships between neuron pairs. Furthermore, the block-diagonal form acts as an intensified regularization for the orthogonal transformation matrix.

\textbf{No inference-time computational overhead}. In contrast to ControlNet, the OFT framework incurs zero additional computational cost during inference. During deployment, the orthogonal transformation matrix $\bm{R}$ obtained through learning is directly applied to the pretrained weight matrix $\bm{W}^0$, resulting in a new effective weight matrix, $\thickmuskip=2mu \medmuskip=2mu \bm{W}=\bm{R}\bm{W}^0$. Consequently, the inference performance remains on par with that of the original pretrained model.

\section{Experiments and Results}

\textbf{Experimental protocol}. Our experiments employ Stable Diffusion v1.5~\cite{rombach2022high} as the base text-to-image generation model. To avoid sampling bias, we randomly select images produced by each technique. The experimental procedures for subject-driven image synthesis are modeled after DreamBooth~\cite{ruiz2023dreambooth}, while those for controllable image generation largely adhere to the guidelines of ControlNet~\cite{zhang2023adding} and T2I-Adapter~\cite{mou2023t2i}. For a consistent comparison with LoRA, both OFT and COFT modifications are restricted to the identical layers where LoRA is applied. Additional experimental results and detailed configurations can be found in Appendix~\ref{app:settings}.

\subsection{Subject-driven Generation}

\textbf{Protocol}. We benchmark our method against DreamBooth~\cite{ruiz2023dreambooth} and LoRA~\cite{hulora2022}. All baselines utilize the DreamBooth loss objective. The implementations for DreamBooth and LoRA generally mirror the recommendations of their original works, with hyperparameters chosen for optimal performance. Further results and elaborations are available in Appendix~\ref{app:settings},\ref{app:coft_oft},\ref{app:more_qualitative},\ref{app:failure}.

\textbf{Finetuning stability and convergence}. Our initial analysis focuses on the stability and convergence dynamics during finetuning across DreamBooth, LoRA, OFT, and COFT methods. Figures~\ref{fig:teaser} and \ref{fig:db_convergence} present a visual summary of our findings. Both COFT and OFT demonstrate a stable capacity to finetune the diffusion model. At the 400-iteration mark, DreamBooth and both OFT variants succeed in achieving effective control, while LoRA is unable to reliably maintain subject identity. By 2000 iterations, DreamBooth begins to produce degenerate (collapsed) outputs and LoRA no longer generates the intended yellow shirt, often substituting yellow fur instead. In contrast, OFT and COFT continue to exhibit robust and consistent generative control. These observations affirm the rapid and stable convergence properties of the OFT framework in the context of subject-driven generation. Notably, the robustness to the number of finetuning steps plays a critical role, as it simplifies the tuning process for different subjects. In practice, both OFT and COFT allow for a straightforward setting of a larger iteration count without the need for intensive tuning. With a well-chosen $\epsilon$ parameter in COFT, setting the learning rate and iteration number becomes largely hassle-free.

\textbf{Quantitative comparison}. Employing the protocol described in \cite{ruiz2023dreambooth}, we systematically compare methods in terms of subject fidelity (using DINO~\cite{caron2021emerging}, CLIP-I~\cite{radford2021learning}), text prompt fidelity (CLIP-T~\cite{radford2021learning}), and sample diversity (LPIPS~\cite{zhang2018unreasonable}). Specifically, CLIP-I is the mean pairwise cosine similarity of CLIP embeddings computed between synthesized and reference images; DINO operates analogously, but leverages ViT S/16 DINO features. CLIP-T reflects the average cosine similarity between textual prompts and generated outputs using CLIP embeddings. Diversity is further measured by computing the average LPIPS cosine similarity among images generated for the same subject and prompt. Results presented in Table~\ref{table:dreambooth_final} indicate that both COFT and OFT substantially surpass DreamBooth and LoRA on the DINO and CLIP-I metrics, and deliver slightly superior or competitive results for prompt fidelity and diversity. For all approaches, each pretrained model undergoes 30 independent finetuning runs with varying random seeds to control for stochastic effects. Overall, the data confirm that the OFT framework achieves not only enhanced convergence stability, but also consistently improved final performance outcomes.

\textbf{Qualitative comparison}. To provide clearer insight into the advantages offered by OFT, Figure~\ref{fig:dreambooth_final} displays a selection of randomly chosen samples for subject-driven generation. For consistency across methods, all images are generated from an identical fine-tuned model, ensuring that no method benefits from individually tailored text prompts for optimal outputs. For each technique, we utilize the model that attains the highest validation CLIP score. Examining the outcomes presented in Figure~\ref{fig:dreambooth_final}, it becomes evident that both OFT and COFT are highly effective in retaining the semantic identity of the subject, whereas LoRA frequently fails to do so—an instance being the bowl example, where subject identity is entirely lost. In addition, OFT and COFT exhibit significantly enhanced responsiveness to textual instructions compared to DreamBooth, which, although it can maintain subject identity, frequently does not generate outputs consistent with the provided prompts (as illustrated in the first row of the bowl example). These qualitative findings illustrate that the OFT framework attains superior controllability in generation while simultaneously ensuring faithful subject representation. Furthermore, the OFT method exhibits robustness to the number of training iterations, maintaining strong performance even with extended training, a property not shared by either DreamBooth or LoRA. Additional qualitative results are provided in Appendix~\ref{app:more_qualitative}. A complementary human evaluation reported in Appendix~\ref{app:human} further confirms OFT's leading performance.

\subsection{Controllable Generation}

\textbf{Settings}. For baseline comparisons, we include ControlNet~\cite{zhang2023adding}, T2I-Adapter~\cite{mou2023t2i}, and LoRA~\cite{hulora2022}. The study addresses three demanding controllable generation tasks: Canny edge to image~(C2I) on the COCO dataset~\cite{lin2014microsoft}, segmentation map to image~(S2I) using ADE20K~\cite{zhou2017scene}, and landmark to face~(L2F) on CelebA-HQ~\cite{karras2018progressive,xia2021tedigan}. Every method undergoes fine-tuning on Stable Diffusion~(SD) v1.5 with the corresponding dataset for 20 training epochs. Further qualitative and failure case results can be found in Appendix~\ref{app:more_qualitative},\ref{app:more_control_task},\ref{app:failure}.

\textbf{Convergence}. To assess the convergence rates of ControlNet, T2I-Adapter, LoRA, and COFT on the L2F task, we conducted both quantitative and qualitative analyses. For our quantitative assessment, we measured the mean $\ell_2$ distance between the landmarks predicted by the model and the control face landmarks. Figure~\ref{fig:face_convergence} illustrates the facial landmark error as a function of the number of finetuning epochs. The data reveal that COFT and OFT exhibit notably accelerated convergence compared to the alternatives. Notably, LoRA requires 20 epochs to achieve the same level of performance that our OFT model attains by epoch 8. Importantly, both OFT and COFT maintain a similar quantity of trainable parameters as LoRA (and substantially fewer than ControlNet), while demonstrating superior convergence efficiency over prior methods. This rapid convergence of OFT is further corroborated in Figure~\ref{fig:teaser}, where the example on the right demonstrates OFT’s substantial data efficiency; specifically, OFT reaches effective convergence using merely 5\% of the ADE20K dataset, outperforming both ControlNet and LoRA under similar conditions. For the qualitative analysis, we center our comparison on OFT, COFT, and ControlNet, since ControlNet's landmark error closely approximates those of our models. As seen in Figure~\ref{fig:face_convergence_qual}, both OFT and COFT converge reliably, with generated face poses aligning incrementally with the control landmarks. In contrast, ControlNet exhibits an abrupt emergence of controllability after epoch 8, corresponding to the sudden convergence behavior reported by \cite{zhang2023adding}. This consistently stable convergence renders OFT particularly practical, as its training dynamics are smoother and more predictable, thereby facilitating the identification of optimal hyperparameters for OFT.

\textbf{Quantitative comparison}. To assess the effectiveness of controllable generation, we propose a control consistency metric. This metric functions by extracting the control signal from each generated image and directly measuring its alignment with the original input control. For the C2I scenario, we report IoU and F1 score; for S2I, we assess performance using mean IoU, as well as mean and overall accuracy; and for L2F, we utilize the mean $\ell_2$ distance between control landmarks and the corresponding predicted landmarks. Further explanation of these consistency metrics is available in Appendix~\ref{app:settings}. Each baseline method is evaluated using its optimal hyperparameters. As seen in Table~\ref{table:control_gen}, both OFT and COFT provide more precise and robust control compared to alternative approaches. Notably, adapter-based methods (\emph{e.g.}, T2I-Adapter and ControlNet) tend to converge more slowly and achieve suboptimal final accuracy. When comparing LoRA and ControlNet, LoRA outperforms ControlNet on the S2I task but underperforms on C2I and L2F. Generally, the upper bound for LoRA’s performance remains low, even with extensive hyperparameter tuning. In contrast, our empirical findings indicate that OFT’s performance continues to improve as the number of trainable parameters increases, suggesting that it has not yet reached its maximum capacity. It is worth highlighting that our rigorous quantitative assessment for controllable generation represents a novel contribution, providing precise evaluation of control fidelity for finetuned models—subject to the inherent accuracy limits of the segmentation or detection models used.

\textbf{Qualitative comparison}. In addition to quantitative assessment, we present a qualitative evaluation of OFT and COFT alongside leading methods such as ControlNet, T2I-Adapter, and LoRA. The randomly sampled outputs in Figure~\ref{fig:control_final} illustrate that OFT and COFT deliver superior image fidelity and realism, while also offering precise controllability. For the S2I task, LoRA is unable to produce outputs that adhere to the specified segmentation maps, whereas ControlNet, OFT, and COFT demonstrate robust control over the generation process. Compared with ControlNet, both OFT and COFT achieve higher image quality, richer detail, and more plausible geometric arrangements using significantly fewer parameters. During the C2I task, OFT and COFT generate semantically coherent images based on minimal Canny edge information, outperforming T2I-Adapter and LoRA, which lag notably behind. For L2F, our approach most accurately follows the target facial poses, even in the presence of difficult orientation challenges. Across all three examined control tasks, both OFT and COFT consistently surpass the competing state-of-the-art models, highlighting the efficacy of our OFT framework in the domain of controllable generation. To supplement these findings, Appendix~\ref{app:control_exp} includes a broader range of qualitative samples for the three control scenarios. Furthermore, Appendix~\ref{app:more_control_task} shows that OFT generalizes effectively to additional control tasks, such as dense pose to human body, sketch to image, and depth to image translations.

\section{Concluding Remarks and Open Problems}

Building on the insight that angular relationships between neurons are fundamental in shaping visual semantics, we introduce a straightforward yet powerful finetuning strategy—orthogonal finetuning—for controlling text-to-image diffusion models. Our approach is applied to two specific scenarios: subject-driven generation and controllable generation. When compared to prior techniques, OFT achieves enhanced controllability and greater stability during finetuning, all while utilizing fewer parameters. Crucially, OFT incurs no extra computational cost at inference time, making the resulting models practical for deployment.

Despite its strengths, OFT raises several open questions. First, OFT maintains orthogonality through Cayley parametrization, which relies on computing a matrix inverse. This requirement somewhat constrains the scalability of the approach. While our block diagonal parametrization alleviates this issue to some extent, identifying more efficient, differentiable methods for handling this matrix inverse remains an unresolved challenge. Second, OFT possesses notable advantages for compositionality: the orthogonal matrices derived from separate OFT finetuning tasks can be multiplied to yield another orthogonal matrix. Whether these combined matrices retain the accumulated knowledge from all corresponding downstream tasks is a compelling area for future research. Lastly, OFT's parameter efficiency is closely tied to the block diagonal structure, which can introduce biases and restrict adaptability. Discovering strategies to further enhance parameter efficiency with reduced bias is a significant open question for ongoing work.

\end{document}